\subsubsection{Model II-4}
% =========================================================
% Model V — Force-conjugate gauge (exact SV, V = ∫ τ da)
% =========================================================

\paragraph{Kinematics.}
\[
\boxed{\
\begin{aligned}
\bm u(x)&=x\,a_\varepsilon\,\FrameBasis\;-\;\tfrac12 x^2\,\bm a_\Sigma\;-\;\tfrac16 x^3\,\bm b_\Sigma,\\
\bm\theta(x)&=x\,(a_\vartheta-c_\vartheta)\,\FrameBasis
\;-\;x\,(\FrameBasis\times\bm a_\Sigma)\;-\;\tfrac12 x^2\,(\FrameBasis\times\bm b_\Sigma)\;-\;(\FrameBasis\times\bm b_\Sigma).
\end{aligned}}
\]


\paragraph{Shear strain and stress.}
Define \(\bm\gamma_\perp=(\mathbf{1}-\FrameBasis\otimes\FrameBasis)\,\bm\gamma\) and
\[
\bm C(\bm{r}):=\varkappa\,(\FrameBasis\times\bm{r})\;+\;\alpha_\varphi\,\nabla_\Sigma\varphi(\bm{r})+\bm{W}^{(b)}\,\bm b_\Sigma,
%\qquad \bm M(\bm{r}):=\operatorname{dev}\bm{r}\otimes\bm{r}+\bm{r}\otimes\CentroidLocation.
\]
A direct computation from \(\hat{\bm u}=\bm u+\bm\theta\times\bm{r}+\cdots\) gives
\[
\bm\varepsilon_\gamma(\hat{\bm u})
=\nabla_\Sigma(\FrameBasis\cdot\hat{\bm u})+\partial_x\hat{\bm u}_\perp
=\bm\gamma_\perp+\bm{\alpha}_r+(\nabla_\Sigma\bm\psi)\,\bm{\alpha}^{(b)}+\bm C(\bm{r}),
\]
and the (engineering) shear stress is \(\bm\tau=G\,\bm\varepsilon_\gamma\).




\subsection{Model III (C, linear shear)}

\paragraph{Kinematics.}

\begin{equation*}
\boxed{\ 
\begin{aligned}
\bm u(x) &=
-\,\tfrac{1}{2} x^2\,\bm{a}_{\Section}\;-\;\tfrac16 x^3\,\bm{b}_{\Section}\;+\;x\,a_\varepsilon\,\FrameBasis\\
\bm\theta(x)&=
-\,x\,(\FrameBasis\times\bm a_\Sigma)
+\Big(x-\tfrac{1}{2} x^2\Big)(\FrameBasis\times\bm b_\Sigma)
+\ (a_\vartheta-c_\vartheta)\,x\,\FrameBasis
\end{aligned}\ }
\end{equation*}



\paragraph{Deformations.}

\begin{equation*}
\begin{aligned}
\bm{\gamma}(x)&= -\,x\,\bm{b}_{\Section}\\
\bm{\kappa}(x)&= (a_\vartheta-c_\vartheta)\,\FrameBasis
-\,(\FrameBasis\times\bm{a}_{\Section})\;+\;(1-x)\,(\FrameBasis\times\bm{b}_{\Section})
\end{aligned}
\end{equation*}
Hence \(\bm{\gamma}\equiv\mathbf{0}\) if and only if \(\bm{b}_{\Section}=\mathbf{0}\).

\paragraph{Warping Basis.}
The warping shapes are:
\begin{align*}
&\bm{W}^{(a)}:=-\,\nu\,\big(\operatorname{dev}\bm{r}\otimes\bm{r}\big),\\
&\bm{W}^{(\varepsilon)}:=-\,\nu\,\bm{r},\\
&\bm{W}^{(b)}:=-\,\nu\,\Big(\operatorname{dev}\bm{r}\otimes\bm{r}+\bm{r}\otimes\CentroidLocation\Big)\mathbf{e}_y,\\[2pt]
&\bm\Phi=\big[\ \bm{W}^{(a)}_y,\ \bm{W}^{(a)}_z,
              \ \bm{W}^{(\varepsilon)},
              \ \bm{W}^{(b)}_y,\ \bm{W}^{(b)}_z,
              \ y\FrameBasis,\ z\FrameBasis,
              \ \FrameBasis,
              \ \FrameBasis\psi_y,
              \ \FrameBasis\psi_z,
              \ \FrameBasis\varphi
              \ \big].
\end{align*}

\paragraph{Amplitudes.}
\begin{equation*}
\boxed{
\begin{aligned}
&\alpha^{(a)}_y=a_y,\\
&\alpha^{(a)}_z=a_z,\\ &\alpha^{(\varepsilon)}=a_\varepsilon,\\
&\alpha^{(b)}_y=\gamma_y,\\
&\alpha^{(b)}_z=\gamma_z,\\
&\alpha_{Y}=-\,\gamma_y,\\
&\alpha_{Z}=-\,\gamma_z,\\
&\alpha_{\mathbf{1}}(x)=-\,\tfrac{1}{2} x^2\,(\CentroidLocation\cdot\bm{b}_{\Section}),\\
&\alpha_{\Psi_y}=b_y,\\
&\alpha_{\Psi_z}=b_z,\\
&\alpha_{\varphi}=a_\vartheta+c_\vartheta.
\end{aligned}}
\end{equation*}
With \(\bm{\gamma}(x)=-x\,\bm{b}_{\Section}\), the Poisson block driven by \(\bm{b}_{\Section}\) satisfies
\(\alpha^{(b)}=\bm{\gamma}\) and the linear-axial block \((Y,Z)\) satisfies
\(\alpha_{Y,Z}=-\,\gamma_{y,z}\), eliminating four warping amplitudes directly by \(\bm{\gamma}\).
The remaining amplitudes \(\alpha_{\mathbf{1}},\alpha_{\Psi_y},\alpha_{\Psi_z},\alpha_{\varphi}\) are not eliminated
(their \(x\)-dependences are \(x^2\), constant, constant, constant, respectively).

\paragraph{Verification.}
Vector identities used: \(\ \FrameBasis\cdot\big((\FrameBasis\times\bm s)\times\bm{r}\big)=-\,\bm s\cdot\bm{r}\),\ \ 
\(\FrameBasis\times(\FrameBasis\times\bm s)=-\,\bm s\) for \(\bm s\perp\FrameBasis\).
\begin{itemize}
\item Transverse alignment terms:
\[
\bm u=-\tfrac16 x^3\bm{b}_{\Section} - \tfrac{1}{2} x^2\bm{a}_{\Section} + x a_\varepsilon \FrameBasis
\],
\(\ \bm\theta\times\bm{r}=\big[x(a_\vartheta-c_\vartheta)\big](\FrameBasis\times\bm{r})
+\big(x-\tfrac{1}{2} x^2\big)(\FrameBasis\times\bm{b}_{\Section})\times\bm{r}
- x(\FrameBasis\times\bm{a}_{\Section})\times\bm{r}\).
\item Shear strain:
\(\bm{\gamma}=\bm u'+\FrameBasis\times\bm\theta
= (-\tfrac{1}{2} x^2\bm{b}_{\Section})+\big(+\tfrac{1}{2} x^2\bm{b}_{\Section}- x\bm{b}_{\Section}\big)=-\,x\,\bm{b}_{\Section}\).
\item Axial fiber strain:
\(\FrameBasis\cdot(\bm\theta\times\bm{r})
= \tfrac{1}{2} x^2\,\bm{r}\cdot\bm{b}_{\Section} - x\,\bm{r}\cdot\bm{b}_{\Section}
+ x\,0\) (torsion part),
and the \((Y,Z)\) block adds \(+\,x\,\bm{r}\cdot\bm{b}_{\Section}\), leaving
\(\tfrac{1}{2} x^2\,\bm{r}\cdot\bm{b}_{\Section}\).
Adding \(\alpha_{\mathbf{1}}\) gives \(+\tfrac{1}{2} x^2(\bm{r}-\CentroidLocation)\cdot\bm{b}_{\Section}\).
\item Warping:
Poisson from \(\bm{b}_{\Section}\): \(\alpha^{(b)}=\bm{\gamma}\) yields
\(-\,x\,\nu\big(\operatorname{dev}\bm{r}\otimes\bm{r}+\bm{r}\otimes\CentroidLocation\big)\bm{b}_{\Section}\).
Poisson from \(\bm{a}_{\Section},a_\varepsilon\): \(\alpha^{(a)},\alpha^{(\varepsilon)}\) reproduce the \(x\)-independent Poisson terms.
Axial \(\psi\): \(\alpha_{\Psi_y}=\!b_y,\ \alpha_{\Psi_z}=\!b_z\) give \(\FrameBasis\otimes\bm\psi\,\bm{b}_{\Section}\).
Torsion warping: \(\alpha_\varphi=a_\vartheta+c_\vartheta\) contributes \(\varphi\,\FrameBasis\,(a_\vartheta+c_\vartheta)\);
together with the rigid twist \(\bm\theta\) this matches \((x\,\FrameBasis\times\bm{r})\,a_\vartheta
- (x\,\FrameBasis\times\bm{r})\,c_\vartheta+\varphi\,\FrameBasis\,(a_\vartheta+c_\vartheta)\).
\end{itemize}

\subsection{Model IV (C, Linear shear)}
% Scheme C: Timoshenko-anchored identification with shear from \bm b_\Sigma only
% -----------------------------------------------------------------------------
% Goal: \bm{\gamma} \neq \mathbf{0} iff \bm b_\Sigma \neq \mathbf{0}, maximize elimination of warping
% through algebraic association with \bm{\gamma} and \bm{\kappa}. All \nu-factors reside in shapes.

% Standing definitions
% --------------------
% Axis unit: \FrameBasis; in-plane position: \bm{r}=(y,z); reference point: \CentroidLocation.
% Parameters: \bm a_\Sigma=(a_y,a_z)^\top,\ a_\varepsilon,\ a_\vartheta,\ \bm b_\Sigma=(b_y,b_z)^\top.
% Torsion coupling: c_\vartheta=c_\vartheta(\bm b_\Sigma) (derived from \psi).
% Gauges for \varphi,\Psi_y,\Psi_z are fixed (e.g., zero mean).

\paragraph{Kinematics.}
\begin{equation*}
\boxed{\ 
\begin{aligned}
\bm u(x)&= x\,a_\varepsilon\,\FrameBasis
-\,\tfrac{1}{2} x^2\,\bm a_\Sigma\;-\;\tfrac16 x^3\,\bm b_\Sigma\\
\bm\theta(x)&=
\underbrace{x\,(a_\vartheta-c_\vartheta)\,\FrameBasis}_{\text{torsion}}
\underbrace{-\,x\,(\FrameBasis\times\bm a_\Sigma)}_{\text{pure bending}}
\;+\;\underbrace{\Big(x-\tfrac{1}{2} x^2\Big)\,(\FrameBasis\times\bm b_\Sigma)}_{\text{shear+bending from }\bm b_\Sigma}
\end{aligned}\ }
\end{equation*}

\paragraph{Deformations.}
\begin{equation*}
\begin{aligned}
\bm{\gamma}(x)&= -\,x\,\bm b_\Sigma \\
\bm{\kappa}(x)&= 
(a_\vartheta-c_\vartheta)\,\FrameBasis
-\,(\FrameBasis\times\bm a_\Sigma)\;+\;(1-x)\,(\FrameBasis\times\bm b_\Sigma)\\
\end{aligned}
\end{equation*}

\paragraph{Warping basis.}
\begin{align*}
% &\bm{W}^{(a)}_y:=-\,\nu\,\big(\operatorname{dev}\bm{r}\otimes\bm{r}\big)\mathbf{e}_y,\qquad
%   \bm{W}^{(a)}_z:=-\,\nu\,\big(\operatorname{dev}\bm{r}\otimes\bm{r}\big)\mathbf{e}_z,\\
% &\bm{W}^{(\varepsilon)}:=-\,\nu\,\bm{r},\\
% &\bm{W}^{(b)}_y:=-\,\nu\,\Big(\operatorname{dev}\bm{r}\otimes\bm{r}+\bm{r}\otimes\CentroidLocation\Big)\mathbf{e}_y,\qquad
%   \bm{W}^{(b)}_z:=-\,\nu\,\Big(\operatorname{dev}\bm{r}\otimes\bm{r}+\bm{r}\otimes\CentroidLocation\Big)\mathbf{e}_z,\\
% &Y\,\FrameBasis:=y\,\FrameBasis,\qquad Z\,\FrameBasis:=z\,\FrameBasis,\qquad \mathbf{1}\,\FrameBasis:=1\cdot\FrameBasis,\\
% &\Psi_y(\bm{r})\,\FrameBasis,\quad \Psi_z(\bm{r})\,\FrameBasis\quad(\psi=\Psi_y b_y+\Psi_z b_z),\qquad
%  \varphi(\bm{r})\,\FrameBasis,\\[2pt]
&\bm\Phi=\big[\ \bm{W}^{(a)}_y,\ \bm{W}^{(a)}_z,\ \bm{W}^{(\varepsilon)},\
                 \bm{W}^{(b)}_y,\ \bm{W}^{(b)}_z,\
                 y\FrameBasis,\ z\FrameBasis,\ 
                 \FrameBasis,\
                 \Psi_y\FrameBasis,\ \Psi_z\FrameBasis,\ \varphi\FrameBasis\ \big].
\end{align*}

\paragraph{Amplitudes.}
\begin{equation*}
\boxed{
\begin{aligned}
&\alpha^{(a)}_y=a_y,\\
&\alpha^{(a)}_z=a_z,\\ &\alpha^{(\varepsilon)}=a_\varepsilon,\\
&\alpha^{(b)}_y=\gamma_y,\\ 
&\alpha^{(b)}_z=\gamma_z,\\
&\alpha_{Y}=-\,\gamma_y,\\ 
&\alpha_{Z}=-\,\gamma_z,\\
&\alpha_{\mathbf{1}}(x)=-\,\tfrac{1}{2} x^2\,(\CentroidLocation\cdot\bm b_\Sigma),\\
&\alpha_{\Psi_y}=b_y,\\
&\alpha_{\Psi_z}=b_z,\\
&\alpha_{\varphi}=a_\vartheta+c_\vartheta.
\end{aligned}}
\end{equation*}
With \(\bm{\gamma}(x)=-x\,\bm b_\Sigma\), the Poisson block driven by \(\bm b_\Sigma\) satisfies
\(\alpha^{(b)}=\bm{\gamma}\) and the linear-axial block \((Y,Z)\) satisfies
\(\alpha_{Y,Z}=-\,\gamma_{y,z}\), eliminating four warping amplitudes directly by \(\bm{\gamma}\).
The remaining amplitudes \(\alpha_{\mathbf{1}},\alpha_{\Psi_y},\alpha_{\Psi_z},\alpha_{\varphi}\) are not eliminated
(their \(x\)-dependences are \(x^2\), constant, constant, constant, respectively).


